\chapter{水の政治：島国の生命線}

\section{川のない都市}

シンガポールに初めて来たとき、ある驚きがあった。この都市にはほとんど天然の水資源がないのだ。

熱帯気候で降雨量が豊富な島国にとって、これは矛盾しているように聞こえる。しかし実際、シンガポールの水問題は長年にわたる根本的な課題だ――この島には大きな天然河川がなく、雨水を貯めにくく、独立当初はほぼすべての淡水供給を隣国マレーシアからの輸入に頼っていた。

水は、シンガポールにとって最も根源的で最も永続的な戦略的不安だ。

\section{貯水池：独立後の長年にわたる努力}

シンガポール滞在中、ある朝、島の北西部にある貯水池を訪れた。午前6時頃で、空はまだ完全に明るくなっていなかったが、すでに多くの市民が水辺で朝の運動をしていた――散歩する人、ジョギングする人、ゆったりとした足取りで。

この水域はシンガポールでは相当広い部類に入る。水は澄んでいて、専用の遊歩道が整備されていた。遠くから見ると木製の桟橋のように見えたが、近づくとプラスチック複合素材だと分かった――耐候性と耐久性を備えており、この国がいかなることにも示す実用主義的なスタイルそのものだ。

景色は確かに美しかったが、この穏やかな水面の背後には、数十年にわたる並々ならぬ努力があった。

独立直後、シンガポールには大型の貯水池がほとんどなかった。当時、大量の雨水は降った後そのまま海に流れ込み、無駄になっていた。独立後、シンガポール政府は数十年をかけて計画的に貯水池群を建設し、分散した降雨を集めて蓄積していった。この技術的・財政的な投資は、一般の人々が想像するはるかに大きなものだった。

\section{三つの水源、一つの目標}

今日、シンガポールの水資源は主に三つのルートから供給されている。

\textbf{一つ目、マレーシアからの購入水。} これは歴史的に形成された主要な水源であり、シンガポールにとって長期的に最大の戦略的弱点でもある。両国間の給水協定は何度も波乱を経て、水価格の交渉が外交摩擦を引き起こしたこともある。そのため、マレーシアへの水依存を減らすことは、歴代シンガポール政府の核心的な政策目標の一つであり続けてきた。

\textbf{二つ目、雨水収集。} シンガポールは島全体に広がる集水システムを通じて、降雨をさまざまな貯水池に導いている。マリーナ貯水池はその中で最も代表的な工事の一つだ――市内中心部の雨水流出を収集・浄化し、かつて直接海に流れ込んでいた湾を、淡水貯蔵施設へと変えた。

\textbf{三つ目、海水淡水化。} シンガポールは複数の海水淡水化プラントを建設し、逆浸透技術で海水を飲料水に変換している。このアプローチはエネルギー消費が大きくコストも高いが、四方を海に囲まれた都市国家にとって、代替不可能な戦略的価値を持つ。

さらに、シンガポールは廃水再利用技術にも力を入れており、高度に処理された再生水（NEWater）を貯水池に注入したり、産業用途に直接使用したりして、「廃水」を水循環の中に取り込んでいる。

\section{節水：文化としての倹約}

シンガポールでは、節水は政策上の推奨にとどまらず、日常生活に深く根付いた習慣だ。政府は長年にわたって段階的水道料金、公衆教育、技術的監視などの手段を用いて、社会全体の節水意識を育ててきた。降雨量が豊富な年であっても、節水の基調は変わらない――一滴の雨の背後に、かつてこの島がどれほど水に乏しかったかを覚えている人がいるからだ。

\section{水と主権}

シンガポールにおける水資源の問題は、決して技術的な問題にとどまらない。国家主権、外交関係、国家アイデンティティと深く絡み合っている。

李光耀は生前、マレーシアがシンガポールへの給水を止めた場合、シンガポールはそれを戦争行為とみなすと何度も表明した。これは脅しではなく、自らの脆弱性を冷静に認識した小国が、その生存の一線を表明したものだ。

だからこそ、シンガポールが数十年にわたって水資源に投資し続けてきたことは、インフラ整備というより、生死にかかわる戦略的自救と言うべきだろう。水の自給自足に向けた一歩一歩は、この国が外部への依存を減らし、自らの存続基盤を固める努力の積み重ねだ。

資源が極めて乏しい小島において、水は刃の最も細い縁だ。
